\section{Real number fields}

\subsection{Cauchy reals}

There are several ways of defining Cauchy real numbers, some of which are non-equivalent.
Here we adopt the terminology of calling ``Cantor reals'' what is often referred to as ``classical Cauchy reals'' and plainly ``Cauchy reals'' for what is usually called ``modulated Cauchy reals'' with fixed modulus \(n ↦ n\).
\begin{definition}[Cantor real numbers]
  A sequence of \(qₙ\) rational numbers is a \emph{Cantor sequence} when
  \[ \for{n:ℕ}{\exist{N:ℕ}{\for{k,l≥N}{\abs{qₖ - qₗ} < 2⁻ⁿ}}}\text. \]

  The set of Cantor sequences modulo relation
  \[ x\sim y ≔ \for{n:N}{\exists{N:ℕ}{\for{k,l≥N}{\abs{xₖ-yₗ}<2⁻ⁿ}}} \]
  is the set of \emph{Cantor reals}.
\end{definition}

\begin{definition}[Cauchy real numbers]
  A sequence of \(qₙ\) rational numbers is a \emph{Cauchy sequence} when
  \[\for{n:ℕ}{\for{k,l≥n}{\abs{qₖ - qₗ} < 2⁻ⁿ}}\text. \]

  The set of Cauchy sequences modulo relation
  \[ x\sim y ≔ \for{n:N}{\for{k,l≥n}{\abs{xₖ-yₗ}<2⁻ⁿ}} \]
  is the set of \emph{Cauchy reals}.
\end{definition}
\begin{note}
  Here the usual modulated Cauchy reals are equivalent to our Cauchy reals, despite fixing the modulus of convergence.
  You can read more at~\cite{nlab:cauchy}.
\end{note}

\begin{proposition}
  Every Cauchy real is a Cantor real, but the converse does not hold in general.
\end{proposition}

\subsection{Dedekind reals}

\begin{definition}[Dedekind construction]\label{def:Rd}
  A pair \(\p{L, U} ∈ 𝒫(ℚ)²\) is a \emph{Dedekind cut} when
  \begin{enumerate}
  \item \(L\) is inhabited, lower, and upwards open
  \item \(U\) is inhabited, upper, and downwards open
  \item \(p ∈ L ∧ q ∈ U ⇒ p < q\) for any \(p\), \(q : ℚ\)
  \item \(p < q ⇒ p ∈ L ∨ q ∈ U\) for any \(p\), \(q : ℚ\)
  \end{enumerate}

  The set of \emph{Dedkind reals} is the set of Dedekind cuts.
\end{definition}

\begin{proposition}
  Every Cantor real is a Dedekind real.
\end{proposition}

\subsection{Escardó-Simpson reals}

\begin{definition}
  Let \(M : (ℕ → \Rd) → \Rd\) be the map \(M(x) ≔ Σₙxₙ2⁻ⁿ⁺¹\) (the indexing starts at \(0\)).
  The Escardó-Simpson reals is the subset
  \[ ⋂\set{X ⊆ \Rd}{ℚ ⊆ X ∧\text{\(X\) is closed under \(M\)}} \text. \]
\end{definition}

\begin{proposition}
  Every Cantor real is an Escardó-Simpson real.
\end{proposition}

\subsection{MacNeille reals}

\begin{definition}[MacNeille construction]\label{def:Rm}
  A pair \(\p{L, U} ∈ 𝒫(ℚ)²\) is a \emph{Dedekind-MacNeille cut} when
  \begin{enumerate}
  \item \(L\) is inhabited, lower, and upwards open
  \item \(U\) is inhabited, upper, and downwards open
  \item[a.] \label{real-m:L} \(p ∈ L ⇔ \exist{q>p}{q ∉ U}\)
  \item[b.] \label{real-m:U} \(q ∈ U ⇔ \exist{p<q}{p ∉ L}\)
  \end{enumerate}

  The set of \emph{MacNeille reals} is the set of Dedekind-MacNeille cuts.
\end{definition}

\begin{proposition}
  Every Dedekind real is a MacNeille real.
\end{proposition}

\subsection{Orders}

We can define the various order relations on reals by defining an order for the MacNeille reals.
Then the orderings of all other order relations come from the restriction of this one.
\begin{definition}
  For a rational \(q : ℚ\) and \(x : \Rm\) define \(x < q ≔ q ∈ Uₓ\) and \(q < x ≔ q ∈ Lₓ\).
  Then for \(x, y : \Rm\) define \(x < y ≔ \exist{q:ℚ}{x < q < y}\) and \(x ≤ y ≔ L_x ⊆ L_y\).
\end{definition}
\begin{lemma}
  We have for all \(x, y : \Rm\) that \(x ≤ y ⇔ U_x ⊇ U_y\).
\end{lemma}
\begin{proof}
  Suppose \(x ≤ y\). Then for \(q ∈ U_y\) we have \(\exist{p<q}{p ∉ L_y}\) by~\ref{real-m:U}.
  By assumption we have \(p ∉ L_x\), so \(q ∈ U_x\).
  The converse argument is symmetric.
\end{proof}
\begin{note}
  We can rephrase \(x ≤ y\) as \(\for{q:ℚ}{q<x ⇒ q<y}\).
\end{note}

Of special note is the following characterisation of the ordering on the Cauchy reals.
\begin{lemma}
  For \(x, y : \Rc\) we have \(x < y ⇔ \exist{n:ℕ}{xₙ + 2⁻ⁿ < yₙ - 2⁻ⁿ}\) and \(x ≤ y ⇔ ¬(y < x) ⇔ \for{n:ℕ}{yₙ - 2⁻ⁿ ≤ xₙ + 2⁻ⁿ}\).
\end{lemma}
\begin{proof}
  Let \(x < y\), so \(x < q < y\).
  Then for infinitely many \(n, m\) we must have \(xₙ + 2⁻ⁿ < q < yₘ - 2⁻ᵐ\).
  Thus there is a single \(n\) such that \(xₙ + 2⁻ⁿ < q < yₙ - 2⁻ⁿ\).

  In the converse direction suppose \(xₙ + 2⁻ⁿ < yₙ - 2⁻ⁿ\).
  Then we may find some \(q:ℚ\) strictly between these two rational numbers, so \(xₙ + 2⁻ⁿ < q < yₙ - 2⁻ⁿ\).
  Thus \(x < q < y\) as desired.
\end{proof}

\begin{proposition}
  For \emph{Dedekind} reals \(x, y : \Rd\) we have \(¬(x < y) ⇔ y ≤ x\).
\end{proposition}
\begin{proof}
  TODO: Zakaj samo dedekind?
  Suppose we have \(q ∈ L_x∩U_y\) and \(U_x⊇U_y\).
  Then we have \(q ∈ Uₓ∩Lₓ\) which is a contradiction.

  Conversely, suppose \(¬(x < y)\) and let \(q ∈ Uₓ\).
  By assumption \(¬\exist{q:ℚ}{x<q<y}\), so \(\for{q:ℚ}{¬(x<q<y)}\).
  Using this at \(x<s<q\) (such an \(s\) exists by openness) we deduce \(¬(s<y)\), so \(s ∉ L_y\).
  Now by~\ref{real-m:U} we have \(q ∈ U_y\), completing the proof.
\end{proof}

The above proposition does not hold for MacNeille reals, but we do have
\begin{lemma}
  For \(x,y:\Rm\) we have \(x ≤ y ⇒ ¬(y<x)\), or \(¬(x≤y ∧ y<x)\).
\end{lemma}
\begin{lemma}
  For MacNeille reals we have \(x < y ∨ x = y ⇒ x ≤ y\).
\end{lemma}


\subsection{Apartness}

\begin{definition}
  A relation \(\apart : 𝒫(X×X)\) is an \emph{apartness relation} on \(X\) when it is
  \begin{itemize}
  \item irreflexive: \(\for{x:X}{¬(x \apart x)}\)
  \item symmetric: \(\for{x,y:X}{x \apart y ⇒ y \apart x}\)
  \item comparative: \(\for{x,y,z:X}{x \apart z ⇒ x \apart y ∨ y \apart z}\)
  \end{itemize}

  An apartness relation is \emph{tight} when \(\for{x,y:X}{¬(x \apart y) ⇒ x = y}\).
\end{definition}

\begin{definition}
  For \(x,y:\Rm\) define \(x \apart y ≔ x < y ∨ y < x\).
\end{definition}

\begin{lemma}
  The relation \(\apart\) is a tight apartness relation on \(\Rd\).
\end{lemma}
\begin{proof}
  Since the relation is obviously irreflexive and symmetric, we only show it's comparative.
  Let \(x \apart z\) so \(x < z ∨ z < x\).
  By symmetry we can consider only \(x < z\).
  We aim to show \(x < y ∨ y < z\).

  Suppose \(¬(x \apart y)\).
  Then \(¬(x < y) ∧ ¬(y < x)\) so \(x ≤ y ∧ y ≤ x\).
  By anti-symmetry we get \(x = y\).
\end{proof}

But \(\apart\) is in general not comparative on the whole of the MacNeille reals.
\begin{theorem}
  In the sheaf topos over \(ℝ\) the relation \(\apart\) on \(\Rm\) is not comparative.
\end{theorem}
\begin{proof}
  Consider \(y ≔ ⋁\set{x:\Rm}{x ≤ 0 ∨ (x ≤ 1 ∧ U)}\), where \(U\) is the generic proposition \(\p{0, ∞}\).
  Now let \(x = 0\) and \(z = 1\).
  Then comparativeness would require \(0 \apart y (≡ U)\) and \(y \apart 1 (≡ ¬U)\) to cover, however \(U ∪ ¬U = ¬\{0\} ≠ ⊤\).
\end{proof}
This proof actually extends to show that there can be no apartness relation on \(\Rm\).
Suppose \(\apart\) is some apartness relation on \(\Rm\).
Then as constructed above \(0 \apart y\) must be at least \(U\), since restricting to \(U\) we have \(y = 1\).
Further, on \(¬U\) we have \(y = 0\), so \(0 \apart y\) is at most \(\int\cl U\).
But because \(U\) is regular \(0 \apart y\) is just \(U\).
Similarly we deduce that \(y \apart 1\) is \(¬U\), which concludes the proof.


%%% Local Variables:
%%% mode: latex
%%% TeX-master: "../main"
%%% End:
